\documentclass[12pt]{article}
\usepackage[utf8]{inputenc}
\usepackage[T1]{fontenc}
\usepackage[margin=1in]{geometry}
\usepackage{setspace}
\usepackage{xcolor}
\usepackage[colorlinks=true, urlcolor=blue!60!black, linkcolor=blue!60!black]{hyperref}
\setlength{\parindent}{0em}
\setlength{\parskip}{0.45em}
\renewcommand{\baselinestretch}{1.0}
\begin{document}
\begin{center}
{\Large \textit{The Prize Cases}, 67 U.S. (2 Black) 635 (1863)\par}
\vspace{0.3em}
{\normalsize Case Brief --- One-Page Edition\par}
\end{center}

\textbf{Memory Jogger.} Lincoln's blockade: the President may repel attack or rebellion without a declaration of war (\href{https://supreme.justia.com/cases/federal/us/67/635/}{67 U.S. 635}).

\textbf{Facts.} After Fort Sumter, with Congress in recess, Lincoln blockaded Southern ports in April 1861. Before Congress recognized the insurrection on July 13, Union ships captured four merchant vessels as prizes; Congress later ratified Lincoln's acts.

\textbf{Procedural History.} District courts condemned the vessels; the claimants appealed directly.

\textbf{Votes.} Affirmed, 5--4. Grier wrote for the Court; Nelson dissented with Taney, Catron, and Clifford.

\textbf{Holding.} The President cannot declare war, but he is authorized and bound to resist invasion or armed rebellion by force, including blockade; courts defer to his judgment that war exists.

\textbf{Reasoning.} Civil war is a fact, ``never solemnly declared,'' and the President had to meet it as it came; property of residents of hostile territory was capturable. Nelson held that only Congress can create a legal state of war, so all captures before July 13 were void.

\textbf{Examples.} \emph{Same-Side:} repelling a surprise attack on a naval base. \emph{Opposite-Side:} a preventive invasion of a nation that has not attacked. \emph{Fence-Sitter:} strikes after a crippling cyberattack.

\textbf{Critique.} Originalists read a narrow defensive power; Louis Fisher and others say presidents overread the case as a blank check.

\textbf{Key Quotation.} ``He has no power to initiate or declare a war either against a foreign nation or a domestic State. \ldots If a war be made by invasion of a foreign nation, the President is not only authorized but bound to resist force by force. He does not initiate the war, but is bound to accept the challenge without waiting for any special legislative authority.'' --- Justice Grier, Opinion of the Court, 67 U.S. at 668.

\end{document}
